% evaluation.tex — 九、模型的评价、改进与推广（2026-08-15, 四件套）

\subsection{9.1 模型的优点}

（1）问题一的分析链条完整。分布规律从总量、月度、周内、日内四个尺度刻画，相互
关系先以散点图检验形态，再选用斯皮尔曼秩相关并配合显著性检验筛选，方法选择有
依据，结论可直接支撑后续建模（周内规律用于问题二预测，长尾分布用于问题三选品）。

（2）问题二贴合实际经营逻辑。目标函数把损耗品打折回收纳入收益口径，预测方法
利用数据中真实存在的周内规律，方案 A 与方案 B 的双方案设计同时给出了"成本加成
惯例下的可行方案"与"定价优化的利润上限"，两个数字之间的差距本身就是结论。

（3）问题三的建模层次清楚。选品与优化分离：先用熵权法加 TOPSIS 的五指标评价
控制候选质量，再以 0-1 规划决定最终方案，并通过品类需求满足率下限直接回应题面
"尽量满足市场需求"的要求，利润与均衡的权衡过程可解释、可复现。

（4）检验充分。26 项一致性测试、历史同期对比与三类参数扰动共同支撑结论的可信
性，检验对象与指标明确，并非"模型较好"式的空泛陈述。

\subsection{9.2 模型的缺点}

（1）茄类与辣椒类三次拟合的 $R^2$ 低于 0.1，价格对这两个品类销量的解释力有限，
其需求估计的置信度不高；（2）销量与批发价预测只利用了星期效应，未建模天气、
节假日与促销等外生因素，异常日的预测偏差无法消除；（3）单品需求采用线性近似，
样本少于 5 天的单品依赖份额缩放兜底，属于近似处理；（4）损耗率与加成率取历史
水平，未随季节与竞争环境动态调整。

\subsection{9.3 模型的改进方向}

针对上述缺点，改进路径与问题四的数据建议对应：引入天气与客流数据修正需求预测
（式\eqref{eq:q4_weather}），引入废弃量数据修正利润口径（式\eqref{eq:q4_waste}），
引入竞争对手价格实现加成率动态化（式\eqref{eq:q4_comp}）。此外，可用分段或
非参数方法替代线性单品需求曲线，并利用每周后验销量对预测参数做滚动更新，使模型
在使用中自我修正。

\subsection{9.4 模型的推广}

模型的三个模块可以独立迁移：同星期预测适用于具有周内规律的零售品类；熵权
TOPSIS 评价适用于任何多指标选品场景；含打折回收项的利润结构适用于一切存在
清仓机制的易损商品。将批发价预测模块替换为对应供应链的价格模型后，本框架可推广
到水果、水产、烘焙等生鲜品类的补货与定价决策，以及多门店、多层级配送的扩展场景。
